\subsection{জেনেটিক ইঞ্জিনিয়ারিং}

\begin{learningbox}
এই অংশ শেষে শিক্ষার্থী পারবে:
\begin{itemize}
\item জেনেটিক ইঞ্জিনিয়ারিংয়ের ধারণা ব্যাখ্যা করতে।
\item DNA, gene, recombinant DNA ও GMO-এর সম্পর্ক বুঝতে।
\item চিকিৎসা, কৃষি ও শিল্পে জেনেটিক ইঞ্জিনিয়ারিংয়ের ব্যবহার লিখতে।
\item নৈতিকতা ও ঝুঁকি আলোচনা করতে।
\end{itemize}
\end{learningbox}

\subsubsection{জেনেটিক ইঞ্জিনিয়ারিংয়ের ধারণা}

জীবের বৈশিষ্ট্য gene দ্বারা নিয়ন্ত্রিত হয়। কোনো জীবের gene পরিবর্তন, সংযোজন বা অপসারণ করে কাঙ্ক্ষিত বৈশিষ্ট্য তৈরি করার প্রযুক্তিকে জেনেটিক ইঞ্জিনিয়ারিং বলা হয়। Genome.gov ও Britannica-এর ব্যাখ্যায় genetic engineering মূলত organism-এর genetic material পরিবর্তনের প্রযুক্তি।

\begin{definitionbox}{সংজ্ঞা}
জেনেটিক ইঞ্জিনিয়ারিং হলো biotechnology-এর একটি শাখা, যেখানে কোনো জীবের DNA বা gene কৃত্রিমভাবে পরিবর্তন করে নতুন বা উন্নত বৈশিষ্ট্য সৃষ্টি করা হয়।
\end{definitionbox}

\subsubsection{ডিএনএ ও জিন}

DNA হলো জীবের genetic instruction বহনকারী molecule। DNA-এর নির্দিষ্ট অংশ, যা কোনো বৈশিষ্ট্য বা protein তৈরির নির্দেশ বহন করে, তাকে gene বলা হয়।

\begin{keypointbox}{মূল ধারণা}
\begin{itemize}
\item \textbf{DNA}: genetic information বহন করে।
\item \textbf{Gene}: DNA-এর কার্যকর অংশ।
\item \textbf{Chromosome}: DNA ও protein দিয়ে গঠিত structure।
\item \textbf{Mutation}: gene বা DNA sequence-এর পরিবর্তন।
\end{itemize}
\end{keypointbox}

\subsubsection{রিকম্বিন্যান্ট ডিএনএ}

Recombinant DNA হলো দুটি ভিন্ন উৎসের DNA অংশ যুক্ত করে তৈরি DNA molecule। যেমন কোনো bacteria-তে মানুষের insulin তৈরির gene প্রবেশ করিয়ে insulin উৎপাদন করা যায়।

\begin{examplebox}{Insulin উৎপাদন}
আগে insulin প্রাণীর pancreas থেকে নেওয়া হতো। recombinant DNA technology ব্যবহার করে bacteria বা yeast-এর মাধ্যমে human insulin উৎপাদন করা সম্ভব হয়েছে, যা diabetes চিকিৎসায় গুরুত্বপূর্ণ।
\end{examplebox}

\subsubsection{জেনেটিক্যালি মডিফাইড অর্গানিজম}

GMO বা Genetically Modified Organism হলো এমন জীব যার genetic material প্রযুক্তির মাধ্যমে পরিবর্তন করা হয়েছে। ফসলের ক্ষেত্রে pest resistance, herbicide tolerance, nutrition improvement বা shelf life বাড়ানোর জন্য GMO তৈরি হতে পারে।

\subsubsection{চিকিৎসায় ব্যবহার}

\begin{itemize}
\item insulin, vaccine, hormone ও therapeutic protein উৎপাদন।
\item gene therapy research।
\item genetic disease diagnosis।
\item personalized medicine।
\item vaccine development ও disease research।
\end{itemize}

\subsubsection{কৃষিতে ব্যবহার}

\begin{itemize}
\item pest-resistant crop।
\item disease-resistant plant।
\item drought বা salinity tolerant variety।
\item nutrient-enriched crop।
\item দ্রুত ও নির্ভুল plant breeding research।
\end{itemize}

\subsubsection{নৈতিকতা ও ঝুঁকি}

জেনেটিক ইঞ্জিনিয়ারিং শক্তিশালী প্রযুক্তি, তাই নৈতিকতা গুরুত্বপূর্ণ। খাদ্য নিরাপত্তা, biodiversity, unintended effect, patent, farmer dependency, human gene editing এবং informed consent নিয়ে সতর্কতা দরকার।

\begin{mistakebox}{সতর্কতা}
জেনেটিক ইঞ্জিনিয়ারিং মানেই সবসময় ক্ষতিকর নয়, আবার সবসময় ঝুঁকিমুক্তও নয়। scientific evaluation, regulation, ethical review ও public awareness জরুরি।
\end{mistakebox}

\begin{boardbox}{বোর্ড ফোকাস}
জেনেটিক ইঞ্জিনিয়ারিংয়ের উত্তরে DNA, gene, recombinant DNA, GMO, insulin production এবং কৃষি ব্যবহার—এই pointগুলো রাখলে উত্তর পূর্ণ হয়।
\end{boardbox}

\begin{practicebox}
\textbf{MCQ}
\begin{enumerate}
\item জীবের বৈশিষ্ট্য নিয়ন্ত্রণকারী DNA-এর অংশ কোনটি?\\
ক. Gene \quad খ. Monitor \quad গ. Router \quad ঘ. Sensor
\item ভিন্ন উৎসের DNA যুক্ত করে তৈরি DNA কোনটি?\\
ক. Recombinant DNA \quad খ. HTML \quad গ. CSS \quad ঘ. LAN
\end{enumerate}

\textbf{সংক্ষিপ্ত প্রশ্ন}
\begin{enumerate}
\item GMO কী?
\item জেনেটিক ইঞ্জিনিয়ারিংয়ের দুটি চিকিৎসা ব্যবহার লেখ।
\end{enumerate}
\end{practicebox}

